% !TeX root = ../regression_logistique.tex
\chapter{Corrigés}
\label{ann:corriges}

\orient{%
comparer une tentative au code de référence et aux réponses attendues}{%
la tâche correspondante a été tentée}{%
à consulter après avoir écrit puis exécuté sa propre version}

Cette annexe porte le code des douze tâches du cahier pratique et les réponses
aux exercices du fil principal. Chaque section porte le nom du fichier ou le
numéro de l'exercice qu'elle corrige.

\section{Cahier pratique}

\subsection*{\texttt{step0\_prepare.py} --- préparation des données}

\codefile[firstline=13,lastline=39]{scripts/tuto/step0_prepare.py}

\textbf{Transfert.} Avec $40\,\%$ de valeurs absentes, la médiane reste une
statistique correcte de la colonne observée, mais l'écart-type calculé après
imputation s'effondre : $40\,\%$ des lignes reçoivent exactement la même valeur.
Le nuage montrerait une barre verticale dense à l'abscisse de la médiane, et le
coefficient de cette variable serait tiré par un groupe artificiel.

\subsection*{\texttt{step1\_score.py} --- le score}

\codefile[firstline=4,lastline=8]{scripts/tuto/step1_score.py}

\textbf{Transfert.} Sans colonne de $1$, le gradient (ch.~\ref{ch:entrainer})
devrait traiter $\partial J/\partial w_0$ à part, avec $x_{i0}$ implicitement
égal à $1$ ; \texttt{one\_step} et \texttt{descend} hériteraient du même cas
particulier, et l'extension à dix matières multiplierait les branches.

\subsection*{\texttt{step2\_sigmoid.py} --- la sigmoïde protégée}

\codefile[firstline=4,lastline=7]{scripts/tuto/step2_sigmoid.py}

\textbf{Transfert.} À $z=-800$, la sigmoïde protégée rend exactement $0{,}0$.
Pour une étiquette valant $1$, la perte directe demanderait $\log 0$ ; la forme
stable de l'étape suivante rend $800$, valeur finie qui laisse la descente
continuer.

\subsection*{\texttt{step3\_single\_cost.py} --- la perte d'une prédiction}

\codefile[firstline=1,lastline=10]{scripts/tuto/step3_single_cost.py}

\textbf{Transfert.} Pour $z\to-\infty$ et $y=1$,
$\operatorname{softplus}(z)-z\to-z=|z|$ : la perte croît linéairement. Une
croissance quadratique donnerait un gradient non borné ; une croissance bornée
rendrait la correction insensible aux prédictions les plus fausses. La dérivée
$p-y$ reste ici comprise entre $-1$ et $1$ (\annexeref{ann:gradient}).

\subsection*{\texttt{step4\_total\_cost.py} --- la perte moyenne}

\codefile[firstline=6,lastline=10]{scripts/tuto/step4_total_cost.py}

\textbf{Transfert.} Un fichier deux fois plus long, de même composition, double
la somme et laisse la moyenne inchangée. Seule la moyenne compare les deux
fichiers. Le pas $\alpha$ multiplie le gradient : avec une somme, la longueur du
déplacement croîtrait avec l'effectif, et un $\alpha$ réglé sur $1\,600$ lignes
ferait diverger la descente sur $3\,200$.

\subsection*{\texttt{step5\_gradient.py} --- gradient et contrôle numérique}

\codefile[firstline=7,lastline=21]{scripts/tuto/step5_gradient.py}

\textbf{Transfert.} En $(0,0,0)$, tous les scores valent zéro et
$p_i-y_i=\pm0{,}5$ ; un facteur $1/m$ omis multiplie les trois composantes par
$m$, ce qu'une comparaison à la mesure détecte --- mais la mesure elle-même
emploie \texttt{total\_cost}, affecté du même facteur, et les deux erreurs se
compensent. Un point d'évaluation à composantes distinctes, comme
$(0{,}3\,;-0{,}7\,;0{,}4)$, sépare les trois colonnes et révèle un indice
décalé, que $(0,0,0)$ laisse invisible puisque les colonnes y jouent des rôles
symétriques.

\subsection*{\texttt{step6\_one\_step.py} --- un pas}

\codefile[firstline=4,lastline=9]{scripts/tuto/step6_one_step.py}

\textbf{Transfert.} Les \Gryf occupent la région $x_1<0$, $x_2>0$
(figure~\ref{fig:nuage}). Au premier tour, $p_i-y_i$ vaut $-0{,}5$ sur eux et
$+0{,}5$ ailleurs ; la composante $\partial J/\partial w_1$ est donc dominée par
des produits $(-0{,}5)\times(\text{$x_1$ négatif})$, positifs, d'où une pente
positive et un $w_1$ qui diminue. Le raisonnement symétrique donne $w_2$
croissant. La trace du tableau~\ref{tab:trace} confirme : $-0{,}2289$ et
$+0{,}1924$.

\subsection*{\texttt{step7\_loop.py} --- la boucle et ses deux arrêts}

\codefile[firstline=4,lastline=17]{scripts/tuto/step7_loop.py}

\textbf{Transfert.} Un critère portant sur les décisions comparerait le vecteur
des prédictions à celui du tour précédent et s'arrêterait après $k$ tours sans
changement. Il coûte un passage supplémentaire sur les lignes pour former les
prédictions, plus le stockage du vecteur précédent. Sur ces données il se
déclencherait quelques centaines de tours avant le critère de perte, au prix
d'une frontière moins nette : les probabilités resteraient plus proches de
$0{,}5$.

\subsection*{\texttt{step8\_four\_houses.py} --- quatre descentes et l'argmax}

\codefile[firstline=12,lastline=22]{scripts/tuto/step8_four_houses.py}

\textbf{Transfert.} Cinq classes demandent cinq descentes et
$5\times(n+1)$ coefficients, soit $55$ pour dix matières. Les statistiques de
préparation restent au nombre de $3n$. Aucune ligne de \texttt{train} ni de
\texttt{predict} ne change : la première itère sur \texttt{HOUSES}, la seconde
prend le maximum de la liste des scores.

\subsection*{\texttt{step9\_holdout.py} --- découpage stratifié}

\codefile[firstline=12,lastline=23]{scripts/tuto/step9_holdout.py}

\codefile[firstline=32,lastline=41]{scripts/tuto/step9_holdout.py}

\textbf{Transfert.} \Slyth, la plus petite classe avec $60$ lignes de test sur
$300$ attendues, subit la variation relative la plus large : un tirage global
laisse son effectif de test fluctuer d'environ $\pm7$ lignes d'une graine à
l'autre. L'intervalle de Wilson global change peu --- il dépend de $n=320$, non
de la composition --- mais les rappels par classe deviennent instables, celui de
\Slyth pouvant varier de plusieurs points sans que le modèle ait changé.

\subsection*{\texttt{step10\_evaluation.py} --- matrice et incertitude}

\codefile[firstline=14,lastline=30]{scripts/tuto/step10_evaluation.py}

\textbf{Transfert.} Rappel $0{,}55$ pour précision $0{,}98$ : le modèle ne
désigne cette classe qu'à coup sûr et en manque près de la moitié. La ligne de
la classe dans la matrice le révèle --- ses effectifs sont dispersés sur les
autres colonnes --- alors que sa colonne reste presque vide hors diagonale. En
binaire, abaisser le seuil rééquilibrerait les deux mesures ; sous la règle du
maximum, il faudrait agir sur le biais du modèle concerné, ce qui déplace la
limite avec les trois autres régions à la fois.

\subsection*{\texttt{transfert\_micro.py} --- douze pièces}

\codefile[firstline=23,lastline=54]{scripts/tuto/transfert_micro.py}

\textbf{Transfert.} Les douze pièces étant séparables, il existe un $w$ qui
classe tout correctement ; multiplier ce $w$ par $\lambda>1$ laisse la frontière
en place et rapproche chaque probabilité de son étiquette, donc fait décroître la
perte. Aucun minimum fini n'existe, la norme croît sans borne
(\annexeref{ann:norme}) et la baisse relative reste au-dessus de $10^{-6}$ :
seul le plafond arrête la boucle. Un plafond dix fois plus grand laisserait la
frontière à la même position et ne ferait que resserrer la transition, comme au
chapitre~\ref{ch:descente}. La pièce de $10{,}20$~mm et $133$~HV donne
$z=+2{,}775>0$ : rebut.

\section{Exercices du fil principal}

\subsection*{Exercice~4.1 --- vecteur normal et distance}

\begin{align*}
\|w\| &= \sqrt{3{,}454^2+3{,}469^2}=4{,}895,\\
-\frac{w_0}{\|w\|} &= \frac{4{,}745}{4{,}895}=0{,}969.
\end{align*}
La distance signée de l'origine à la frontière vaut donc $0{,}969$ écart-type.
Le score de $x=(0\,;0)$ vaut $w_0=-4{,}745$ : l'observation moyenne appartient
au demi-plan des trois autres classes. Les \Gryf occupent le quadrant
$x_1<0$, $x_2>0$, éloigné du centre.

\subsection*{Exercice~5.1 --- probabilité de bout en bout}

$x_1=(4{,}12-1{,}189)/5{,}175=+0{,}566$. La note d'\Runes étant absente, elle
reçoit la médiane $463{,}918$, d'où
$x_2=(463{,}918-495{,}052)/105{,}186=-0{,}296$. Alors
\[
  z=-4{,}745+(-3{,}454)(0{,}566)+3{,}469(-0{,}296)=-7{,}724,
  \qquad p=0{,}00044 .
\]
Le score est négatif : l'observation appartient au demi-plan des autres classes,
loin de la frontière. L'écart avec l'élève~$3$ ($z=+1{,}34$) vient d'\Herb : le terme
passe de $+5{,}13$ à $-1{,}96$, soit sept unités, quand celui d'\Runes n'en
change que d'une seule.

\subsection*{Exercice~7.1 --- signe d'une composante du gradient}

Voir le transfert de \texttt{step6\_one\_step.py} ci-dessus : $\partial
J/\partial w_1$ est positive et $\partial J/\partial w_2$ négative, d'où un $w_1$
décroissant et un $w_2$ croissant dès le premier tour.

\subsection*{Exercice~8.1 --- lecture d'un coefficient}

Un modèle ajusté sur les notes brutes vérifie
$w_2^{\text{brut}}\,v = w_2^{\text{std}}\,(v-\mu)/\sigma$ à un terme constant
près, absorbé par le biais : $w_2^{\text{brut}}=w_2^{\text{std}}/\sigma
=3{,}469/105{,}186=0{,}033$. Les deux décrivent la même frontière. Seul
$w_2^{\text{std}}$ se compare à celui d'\Herb, les deux étant exprimés par
écart-type. Un écart-type de plus en \Runes multiplie la cote par
$e^{3{,}469}=32$ dans les deux paramétrages : l'exposant est le produit
coefficient $\times$ variation, invariant par changement d'échelle.

\subsection*{Exercice~9.1 --- portée des quatre sorties}

$0{,}952$ répond à « cette observation appartient-elle à \Gryf, oui ou non ? » et
$0{,}899$ à « appartient-elle à \Rav ? ». Les deux questions sont séparées et
leurs réponses ne se partagent aucun budget : leur somme, $1{,}851$ pour cet
élève, dépasse $1$. Seule la seconde phrase est autorisée. La première supposerait
une distribution sur les quatre classes, c'est-à-dire une somme égale à $1$, que
la construction un-contre-tous n'impose pas.

\subsection*{Exercice~I.1 --- pas critique d'une quadratique}

Pour $f(\theta)=\tfrac{c}{2}(\theta-2)^2$, $f'(\theta)=c(\theta-2)$ et
l'itération $\theta_{t+1}=\theta_t-\alpha c(\theta_t-2)$ donne
\[
  \theta_{t+1}-2=(1-\alpha c)\,(\theta_t-2).
\]
La suite converge si $|1-\alpha c|<1$, soit $0<\alpha<2/c$ ; elle oscille sans
converger en $\alpha=2/c$ et diverge au-delà. Avec $c=4$, le pas critique vaut
$0{,}5$. La figure~\ref{fig:ann-opt-alpha-quadratique} montre les deux régimes.
